\ch{State of the Art}{tippsundtricks}
 
This chapter traces the methodological evolution of robotic control toward language-driven autonomy. \textbf{Section 2.1} examines classical robot programming, outlining the transition from vendor-specific control languages to interoperable middleware. \textbf{Section 2.2} extends this by introducing task-planning frameworks that link symbolic reasoning with geometric feasibility. Both streams converge in \textbf{Section 2.3}, where AI-driven robot autonomy is discussed as the shift from deterministic, rule-based control to adaptive, data-driven systems capable of self-directed behavior. Building on this, \textbf{Section 2.4} explores NL-facilitated human-robot cooperation, establishing linguistic interaction as a scalable interface for task specification. This leads to \textbf{Section 2.5}, which formalises LLMs and function calling as structured reasoning mechanisms for transforming NL goals into deterministic tool invocations. \textbf{Section 2.6} integrates these capabilities within robotic architectures, analysing how LLMs can perform high-level planning while maintaining operational safety. Finally, \textbf{Section 2.7} surveys existing LLM-guided applications in industrial and laboratory environments, identifying persistent challenges in grounding, reproducibility, and local deployment.

\includefile{Kapitel_2}{classicalprogramming}
\includefile{Kapitel_2}{task}
\includefile{Kapitel_2}{aidriven}
\includefile{Kapitel_2}{language}
\includefile{Kapitel_2}{functioncalls}
\includefile{Kapitel_2}{llminrobotics}
%\includefile{Kapitel_2}{llmintegration}
\includefile{Kapitel_2}{roboticapplications}
%\includefile{Kapitel_2}{workflow}
\includefile{Kapitel_2}{summary}